% 03 --- Molecular Representation on Computers (rigorous)
\section{Molecular Representation on Computers}

\subsection{Coordinates, Topology}
\begin{frame}{Atomic Coordinates --- The $3N$ Truth}
  \begin{equation*}
    \mathbf{R} = \{\mathbf{r}_1,\dots,\mathbf{r}_N\},\quad \mathbf{r}_i=(x_i,y_i,z_i)\in\mathbb{R}^3
  \end{equation*}
  \begin{columns}[T,onlytextwidth]
    \column{0.5\textwidth}
    \begin{pitemize}
      \item $3N$ numbers = \textbf{complete} classical description
      \item File: \inlinecode{XYZ / PDB / GRO} \ensuremath{\rightarrow} list
      \item Units: Å (most), nm (GROMACS)
      \item \textbf{Rigorous}: coordinates \textbf{are} the state
    \end{pitemize}
    \column{0.5\textwidth}
    \begin{center}
      \begin{tikzpicture}[scale=0.9]
        \node[draw=accentcyan,fill=darkgray,rounded corners=1mm,minimum width=2.6cm,minimum height=0.8cm] at (0,1) {\color{fgwhite}\tiny\ttfamily ATOM 1 C 0.0 0.0 0.0};
        \node[draw=accentcyan,fill=darkgray,rounded corners=1mm,minimum width=2.6cm,minimum height=0.8cm] at (0,0) {\color{fgwhite}\tiny\ttfamily ATOM 2 O 1.2 0.0 0.0};
        \node[draw=accentcyan,fill=darkgray,rounded corners=1mm,minimum width=2.6cm,minimum height=0.8cm] at (0,-1) {\color{fgwhite}\tiny\ttfamily ATOM 3 H -0.4 0.9 0.0};
        \draw[->,thick,neongreen] (1.4,1) -- (2.2,0.2) node[right,color=neongreen]{array};
      \end{tikzpicture}
    \end{center}
  \end{columns}
\end{frame}

\begin{frame}{Bonds, Angles, Dihedrals --- Topology}
  \begin{columns}[T,onlytextwidth]
    \column{0.5\textwidth}
    \begin{pitemize}
      \item \textbf{Bond}: $r_{ij}=|\mathbf{r}_j-\mathbf{r}_i|$ · 2 atoms
      \item \textbf{Angle}: $\theta_{ijk}$ · 3 atoms
      \item \textbf{Dihedral}: $\phi_{ijkl}$ · 4 atoms, \textbf{periodic}
      \item Topology = \textbf{graph} (bonds = edges)
    \end{pitemize}
    \column{0.5\textwidth}
    \begin{tikzpicture}[scale=0.9]
      \coordinate (A) at (0,0); \coordinate (B) at (1.2,0.6); \coordinate (C) at (2.4,0.2); \coordinate (D) at (3.2,1.1);
      \foreach \p in {A,B,C,D} {\shade[ball color=accentcyan] (\p) circle (0.22cm);}
      \draw[thick,white] (A)--(B)--(C)--(D);
      \draw[dashed,neonyellow] (A)--(C); \node[color=neonyellow] at (1.2,-0.4) {angle $\theta$};
      \draw[dashed,neonpink] (B)--(D); \node[color=neonpink] at (2.8,0.7) {dihedral $\phi$};
      \only<2>{\draw[->,thick,neongreen] (B) -- ++(0.4,0.8) node[above]{4-body!};}
    \end{tikzpicture}
  \end{columns}
\end{frame}

\begin{frame}{Interpretation --- Coordinates/Topology}
  \interpret{
    \textbf{coordinates} · $3N$ numbers = molecule · \textbf{bond} · distance · \textbf{angle} · 3 atoms · \textbf{dihedral} · 4 atoms = flexibility · \textbf{topology} · graph
  }
\end{frame}

\subsection{Masses, Charges, Box, PBC}
\begin{frame}{Masses \& Partial Charges}
  \begin{columns}[T,onlytextwidth]
    \column{0.5\textwidth}
    \begin{pitemize}
      \item $m_i$ \ensuremath{\rightarrow} dynamics: $\mathbf{F}=m\mathbf{a}$
      \item $q_i$ \ensuremath{\rightarrow} electrostatics: $U_{\text{Coul}}=\frac{q_iq_j}{4\pi\epsilon_0 r_{ij}}$
      \item \textbf{Partial} \ensuremath{\neq} formal: fitted to reproduce $ \Phi(\mathbf{r})$
      \item Rigorous: charges are \textbf{parameters}, not observables
    \end{pitemize}
    \column{0.5\textwidth}
    \begin{tikzpicture}[scale=0.9]
      \shade[ball color=neonblue] (0,0) circle (0.45cm) node[color=white]{O$^{-0.8}$};
      \shade[ball color=neonyellow] (-0.9,0.7) circle (0.32cm) node[color=bgblack]{H$^{+0.4}$};
      \shade[ball color=neonyellow] (0.9,0.7) circle (0.32cm) node[color=bgblack]{H$^{+0.4}$};
      \draw[->,thick,accentcyan] (0,-0.6) -- (0,-1.3) node[below,color=accentcyan]{$\mathbf{\mu}$};
    \end{tikzpicture}
  \end{columns}
\end{frame}

\begin{frame}{Simulation Box \& Periodic Boundary Conditions}
  \begin{columns}[T,onlytextwidth]
    \column{0.5\textwidth}
    \begin{pitemize}
      \item Box: $\mathbf{L}=(L_x,L_y,L_z)$, triclinic: $\mathbf{h}$
      \item \textbf{PBC}: infinite tiles, minimum image
      \item Rigorous: removes surface, \textbf{finite-size} remains
      \item Cutoff $r_c < L/2$ 
    \end{pitemize}
    \column{0.5\textwidth}
    \begin{center}
      \begin{tikzpicture}[scale=0.6]
        \draw[accentcyan,thick] (0,0) rectangle (1.8,1.4);
        % PBC images (simplified, no conditional)
        \foreach \dx in {-1.8,1.8} \foreach \dy in {-1.4,0,1.4} {
          \draw[softgray,thin] (\dx,\dy) rectangle ++(1.8,1.4); \fill[neonpink] (\dx+0.9,\dy+0.7) circle (1pt);
        }
        \foreach \dx in {0} \foreach \dy in {-1.4,1.4} {
          \draw[softgray,thin] (\dx,\dy) rectangle ++(1.8,1.4); \fill[neonpink] (\dx+0.9,\dy+0.7) circle (1pt);
        }
        \fill[accentcyan] (0.9,0.7) circle (2pt);
        \node[color=accentcyan] at (0.9,1.7) {central box};
        \draw[dashed,neonyellow] (0.9,0.7) -- (2.7,0.7) node[midway,above]{image};
      \end{tikzpicture}
    \end{center}
  \end{columns}
\end{frame}

\subsection{PES \& Force Fields}
\begin{frame}{Potential Energy Surface --- The Mountain Range}
  \begin{equation*}
    U(\mathbf{R}) : \mathbb{R}^{3N}\to\mathbb{R} \qquad \mathbf{F}_i = -\nabla_{\mathbf{r}_i}U
  \end{equation*}
  \begin{columns}[T,onlytextwidth]
    \column{0.5\textwidth}
    \begin{pitemize}
      \item $3N$-dim surface \ensuremath{\rightarrow} $2$D cartoon is \textbf{projection}
      \item Minima = stable structures
      \item Saddle = transition state
      \item Rigorous: \textbf{all} physics in $U$
    \end{pitemize}
    \column{0.5\textwidth}
    \begin{tikzpicture}[scale=0.6]
      \draw[accentcyan!50,fill=accentcyan!10] plot[smooth] coordinates {(0,0) (0.6,0.8) (1.2,0.3) (1.8,1.1) (2.4,0.6) (3,1.3)} -- (3,0) -- cycle;
      \fill[neongreen] (0.6,0.8) circle (2pt) node[above,color=neongreen]{min};
      \fill[neonpink] (1.8,1.1) circle (2pt) node[above,color=neonpink]{min};
      \fill[neonyellow] (1.2,0.3) circle (2pt) node[below,color=neonyellow]{saddle};
      \draw[->,thick,accentcyan] (1.2,0.3) -- (1.8,1.1) node[midway,sloped,above]{barrier};
    \end{tikzpicture}
  \end{columns}
\end{frame}

\begin{frame}{Force Fields --- $U$ Made Computable}
  \begin{equation*}
    U = \underbrace{U_{\text{bond}}+U_{\text{angle}}+U_{\text{dihed}}}_{\text{bonded}} + \underbrace{U_{\text{LJ}}+U_{\text{Coul}}}_{\text{nonbonded}}
  \end{equation*}
  \begin{columns}[T,onlytextwidth]
    \column{0.5\textwidth}
    \begin{pitemize}
      \item Bond: $k_b(r-r_0)^2$ · spring
      \item LJ: $4\epsilon[(\sigma/r)^{12}-(\sigma/r)^6]$
      \item Many: AMBER, CHARMM, OPLS, ReaxFF, \textbf{ML potentials}
    \end{pitemize}
    \column{0.5\textwidth}
    \begin{infobox}[FF = Expensive Physics Compressed]
      \scriptsize $10^5$ DFT points \ensuremath{\rightarrow} $U_{\text{FF}}$ \ensuremath{\rightarrow} $10^9$ steps
    \end{infobox}
  \end{columns}
\end{frame}

\begin{frame}{Interpretation --- PES/FF}
  \interpret{
    \textbf{PES} · energy landscape · \textbf{force} = slope · \textbf{force field} · cheap $U$ · \textbf{bonded + nonbonded} · \textbf{parameterized} \ensuremath{\rightarrow} not universal
  }
\end{frame}
